\begingroup
\section{Node PT: physical material transport and the mixed row DM}
\label{module:PT}
\noindent\textit{Audited source: \nolinkurl{PT_material_transport.tex}.}
\subsection{The row and the two regimes}

Let
\begin{equation}\label{PT:eq:fan}
 h_i>0,\qquad \sum_i h_i=2\pi,\qquad
 a=\frac{2\pi}{N},\qquad S=a+\max_i h_i,
\end{equation}
and put
\begin{equation}\label{PT:eq:J}
 \Jfour(h)=\sum_i h_i^4-Na^4.
\end{equation}
On $I_i=[\theta_i,\theta_{i+1}]$ define
\begin{equation}\label{PT:eq:q}
 q_h(\theta_i+x)=\frac12\min\{x^2,(h_i-x)^2\},
 \qquad0\le x\le h_i.
\end{equation}

Let $Z_h$ be the exact area tangent space of support increments in a fixed
translation gauge, and let $v_h(z)$ be the physical-affine normal trace.
On
\[
 K_i=[\theta_i-h_{i-1}/2,\theta_i+h_i/2],\qquad x=\theta-\theta_i,
\]
it is
\begin{align}
 v_h(z)(\theta)
 &=\PiNG\left[z_i\cos x+
 \{(1-t_i(x))A_i^-+t_i(x)A_i^+\}\sin x\right],
 \label{PT:eq:trace}\\
 t_i(x)&=\frac{\tan x+\tan(h_{i-1}/2)}
 {\tan(h_{i-1}/2)+\tan(h_i/2)},\notag\\
 A_i^-&=\frac{z_i-z_{i-1}}{\sin h_{i-1}}
       -z_i\tan\frac{h_{i-1}}2,\qquad
 A_i^+=\frac{z_{i+1}-z_i}{\sin h_i}
       +z_i\tan\frac{h_i}2.\notag
\end{align}
The exact support metric satisfies
\begin{equation}\label{PT:eq:metric}
 c\|v_h(z)\|_{H^{1/2}}^2
 \le\Gtilde_h(z)\le C\|v_h(z)\|_{H^{1/2}}^2.
\end{equation}

The symmetric null form is
\begin{equation}\label{PT:eq:null}
 \Cnull(f_1,f_2,f_3)
 =\frac13\sum_{\rm cyc}\int_\T f_1
 \{(|D|f_2)(|D|f_3)-f_2'f_3'\}\,d\theta.
\end{equation}
The row DM is
\begin{equation}\label{PT:eq:DM-target}
 |\Cnull(q_h,q_h,v_h(z))|
 \le\varepsilon_{\rm DM}(S)\Jfour(h)^{1/2}
       \Gtilde_h(z)^{1/2},
 \qquad\varepsilon_{\rm DM}(S)\longrightarrow0.
\end{equation}

Fix $\beta=1/2$.  As in stage 161, split into
\begin{align}
 \Jfour(h)&\ge S^{5-\beta}&&\text{(far)},\label{PT:eq:far-regime}\\
 \Jfour(h)&<S^{5-\beta}&&\text{(near)}.\label{PT:eq:near-regime}
\end{align}
The fully polarized one-Lipschitz estimate
\begin{equation}\label{PT:eq:one-slot}
 |\Cnull(f,g,w)|
 \le C\|f'\|_\infty
       \|g\|_{\dot H^{1/2}}\|w\|_{\dot H^{1/2}}
\end{equation}
follows from the same-sign Fourier formula by putting the low derivative
block in $L^\infty$; when the Lipschitz slot is high, the remaining
separation factor is $2^{-|j-k|/2}$.  Since
\begin{equation}\label{PT:eq:q-size}
 \|q_h'\|_\infty\le S,
 \qquad \|q_h\|_{\dot H^{1/2}}\le CS^{3/2},
\end{equation}
the far regime gives
\begin{equation}\label{PT:eq:far-bound}
 |\Cnull(q_h,q_h,v_h(z))|
 \le CS^{1/4}\Jfour(h)^{1/2}\Gtilde_h(z)^{1/2}.
\end{equation}

In the near regime set
\begin{equation}\label{PT:eq:path}
 d_i=h_i-a,\qquad h_i(t)=a+td_i.
\end{equation}
The exact identity
\begin{equation}\label{PT:eq:Jidentity}
 \Jfour(h)=\sum_i d_i^2(h_i^2+2ah_i+3a^2)
\end{equation}
gives $|d_i|^4\le C\Jfour(h)$.  In the near regime,
$\max_i|d_i|\le CS^{(5-\beta)/4}=CS^{9/8}=o(S)$ as $S\to0$.
Since $S=a+\max_i h_i$ and $|h_i-a|\le\max_j|d_j|$, both $a$ and every
$h_i(t)$ therefore lie between fixed positive multiples of $S$, for all
sufficiently small $S$.  Hence
\begin{equation}\label{PT:eq:comparable}
 cS\le a\le h_i(t)\le CS,
 \qquad
 \delta(t):=\max_i\frac{|d_i|}{h_i(t)}
 \le C\frac{\max_i|d_i|}{S}.
\end{equation}

\subsection{Lipschitz covariance of the same-sign null form}

For a periodic Lipschitz function $u$, put
$F_\varepsilon(x)=x+\varepsilon u(x)$ and
\begin{equation}\label{PT:eq:transport-derivative}
 \mathfrak D_u\Cnull(f,g,w)
 =\frac d{d\varepsilon}\bigg|_{\varepsilon=0}
 \Cnull(f\circ F_\varepsilon^{-1},
        g\circ F_\varepsilon^{-1},
        w\circ F_\varepsilon^{-1}).
\end{equation}

\begin{lemma}[Critical Lipschitz covariance]\label{PT:lem:covariance}
For real $f\in W^{1,\infty}(\T)$ and real
$g,w\in H^{1/2}(\T)$,
\begin{equation}\label{PT:eq:covariance}
 \boxed{
 |\mathfrak D_u\Cnull(f,g,w)|
 \le C\|u'\|_\infty\|f'\|_\infty
       \|g\|_{H^{1/2}}\|w\|_{H^{1/2}}.}
\end{equation}
The polarized estimate is uniform for piecewise-affine $u$.
\end{lemma}

\begin{proof}
First take trigonometric polynomials.  The same-sign representation of
\eqref{PT:eq:null} is a sum, over $p,q>0$, of the three placements of
\begin{equation}\label{PT:eq:same-sign}
 pq\,\widehat f_p\widehat g_q\overline{\widehat w_{p+q}}
\end{equation}
and their conjugates.  Since
\[
 \frac d{d\varepsilon}\bigg|_0
 (f\circ F_\varepsilon^{-1})=-uf',
\]
differentiating the three slots in \eqref{PT:eq:same-sign} produces a
quadrilinear multiplier with the deformation frequency $\ell$.  At
$\ell=0$ the three contributions add to zero: this is exactly invariance
under a common translation.  Therefore the summed multiplier is divisible
by $\ell$, and $\ell\widehat u_\ell$ is a Fourier coefficient of $u'$.

For completeness, this divisibility has a pointwise symbol bound.  Put
\[
 \tau(r,s,t)=\sum_{\rm cyc}(|s||t|+st),
 \qquad r+s+t=0.
\]
Thus $\tau$ is twice the product of the two same-sign frequencies.  If
$\ell+k+m+n=0$, direct substitution in the six possible sign orders gives
\begin{multline}\label{PT:eq:symbol-bound}
 \big|k\tau(k+\ell,m,n)+m\tau(k,m+\ell,n)\\
       +n\tau(k,m,n+\ell)\big|
 \le3|\ell|\,|k|\,|m|^{1/2}|n|^{1/2}.
\end{multline}
For example, on a region with $m,n>0$ the only changes occur when
$k+\ell$, $m+\ell$, or $n+\ell$ crosses zero; on each resulting interval
the left side is affine in $\ell$ after its zero value at $\ell=0$ is
removed, and its two endpoint slopes are bounded by the right side.
The other sign orders are permutations.  This proves
\eqref{PT:eq:symbol-bound} without a smoothness assumption on $u$.

We spell out the endpoint summation, since the pointwise bound alone is
not an $L^\infty\times L^\infty\times L^2\times L^2$ theorem.  Let
$P_j$ be cyclic Littlewood--Paley projections and write
\[
 U_j=P_j(u'),\quad F_j=P_j(f'),\quad
 G_j=|D|^{1/2}P_jg,\quad W_j=|D|^{1/2}P_jw.
\]
Split each dyadic block further at the six sign boundaries
\[
 k=0,\quad m=0,\quad n=0,\quad
 k+\ell=0,\quad m+\ell=0,\quad n+\ell=0.
\]
On every resulting interval the divided symbol
\begin{equation}\label{PT:eq:divided-symbol}
 \frac{k\tau(k+\ell,m,n)+m\tau(k,m+\ell,n)
       +n\tau(k,m,n+\ell)}
      {\ell k|m|^{1/2}|n|^{1/2}}
\end{equation}
and each discrete difference, multiplied by the corresponding dyadic
scale, are bounded.  At a boundary the one-sided limits agree with the
formula in the adjacent sign chamber; if one of $k,m,n$ is zero, the
undivided symbol is zero.  Thus no boundary atom or zero-frequency row is
lost.

Let $J$ be the largest dyadic index and recall that the frequency sum
forces the two largest indices to differ by at most $3$.  Up to exchanging
$g,w$ and $u',f'$, all flag rows have the following bounds:
on the chambers where one of $G,W$ is $r$ generations below a high
$U$ or $F$ block, direct substitution in the same-sign product gives an
extra $2^{-r/2}$; when both $G,W$ are lower by $r,s$ and $U,F$ are high,
the divided symbol is bounded by $C2^{-(r+s)/2}$.  For example, if
$k>0$, $m,n>0$, and $\ell<0$ are in the latter configuration, then
$k+\ell=-(m+n)$ and the numerator in
\eqref{PT:eq:divided-symbol} is $6kmn$, so division gives
$C|mn|^{1/2}/|\ell|$.  The other five sign orders are identical after a
permutation.  Thus all flag rows have the following bounds:
\begin{center}
\begin{tabular}{c|c}
two high blocks & bound after the order-zero kernel is summed\\
\hline
$G_J,W_J$ &
$C\|U_{\le J}\|_\infty\|F_{\le J}\|_\infty
  \|G_J\|_2\|W_{J+O(1)}\|_2$\\
$U_J,G_J$ &
$C\sum_{r\ge0}2^{-r/2}\|U_J\|_\infty
  \|F_{\le J}\|_\infty\|G_{J+O(1)}\|_2\|W_{J-r}\|_2$\\
$F_J,G_J$ &
$C\sum_{r\ge0}2^{-r/2}\|U_{\le J}\|_\infty
  \|F_J\|_\infty\|G_{J+O(1)}\|_2\|W_{J-r}\|_2$\\
$U_J,F_J$ &
$C\sum_{r,s\ge0}2^{-(r+s)/2}\|U_J\|_\infty
  \|F_{J+O(1)}\|_\infty\|G_{J-r}\|_2\|W_{J-s}\|_2$.
\end{tabular}
\end{center}
The factors $2^{-r/2}$ are the unused half derivative in
\eqref{PT:eq:symbol-bound}.  The same table holds when a shifted frequency
touches a sign boundary, using the one-sided divided symbol above.

Each Littlewood--Paley block is bounded in $L^\infty$ by the full
$L^\infty$ norm.  Sum first in $r,s$ and then in $J$.  Young's inequality
for the geometric sequences and Cauchy--Schwarz for the $G,W$ square
functions give
\[
 C\|u'\|_\infty\|f'\|_\infty
 \left(\sum_j\|G_j\|_2^2\right)^{1/2}
 \left(\sum_j\|W_j\|_2^2\right)^{1/2}.
\]
The finitely many zero and unit frequencies are bounded by the
inhomogeneous $H^{1/2}$ norms.  This is \eqref{PT:eq:covariance}.

Approximate a Lipschitz $u$ by smooth functions with uniformly bounded
derivative.  The estimate just proved is stable under this approximation,
so it applies in particular to piecewise-affine velocities.
\end{proof}

\begin{remark}
Lemma~\ref{PT:lem:covariance} measures $u'$ and not $u$.  Consequently a
large cumulative displacement of grid nodes and a common rotation do not
enter the estimate.  This is the precise replacement for the invalid
bare-kernel order used in stage 161.
\end{remark}

\subsection{Exact material fan identity and support transport}

Refine the fan mesh by the nodes $\theta_i(t)$ and the midpoints
$c_i(t)=(\theta_i(t)+\theta_{i+1}(t))/2$.  Let
$F_t:\T_0\to\T_t$ be the unique orientation-preserving map which is
affine on every reference half-cell and sends
\(\theta_i(0),c_i(0)\) to \(\theta_i(t),c_i(t)\).  Set
\begin{equation}\label{PT:eq:finite-material-map}
 G_{t,\varepsilon}=F_{t+\varepsilon}\circ F_t^{-1},
 \qquad
 u_t=\partial_\varepsilon G_{t,\varepsilon}\big|_{\varepsilon=0}.
\end{equation}
On the two halves of $I_i(t)$ one has
\begin{equation}\label{PT:eq:uprime}
 u_t'=\frac{d_i}{h_i(t)},
 \qquad \|u_t'\|_\infty=\delta(t).
\end{equation}

\begin{lemma}[Exact fan transport]\label{PT:lem:q-transport}
On both halves of $I_i(t)$,
\begin{equation}\label{PT:eq:q-transport}
 q_{h(t+\varepsilon)}\circ G_{t,\varepsilon}
 =\left(1+\varepsilon\frac{d_i}{h_i(t)}\right)^2q_{h(t)}.
\end{equation}
Consequently, if a dot denotes the derivative of the pulled profile,
\begin{equation}\label{PT:eq:qdot}
 \|\dot q_t\|_{H^{1/2}}
 \le C\delta(t)S^{3/2}.
\end{equation}
\end{lemma}

\begin{proof}
On either half-cell, $q_h$ is one half of the squared distance from its
nearest endpoint.  The affine map multiplies that distance by
$1+\varepsilon d_i/h_i(t)$, proving \eqref{PT:eq:q-transport}.

Thus $\dot q_t=2(d_i/h_i(t))q_{h(t)}$ on $I_i(t)$.  Each cell piece
vanishes quadratically at its outer endpoints.  Its zero extension has
$H^{1/2}$ norm squared at most $Ch_i(t)^4|d_i/h_i(t)|^2$ by scaling the
fixed unit-cell profile.  To sum the pieces, use comparability from
\eqref{PT:eq:comparable}.  The $H^{1/2}$ pairing of two nonadjacent cell
profiles at cyclic index distance $r$ is bounded by
$CS^4|a_i a_j|/(1+r)^2$, where $a_i=d_i/h_i(t)$; adjacent pairs are
bounded by the two-cell unit-profile norm.  Schur summation of
$(1+r)^{-2}$ gives
\[
 \left\|\sum_i2a_iq_i\right\|_{H^{1/2}}^2
 \le C\sum_i S^4|a_i|^2.
\]
Hence
\[
 \|\dot q_t\|_{H^{1/2}}^2
 \le C\delta(t)^2\sum_i h_i(t)^4
 \le C\delta(t)^2S^3,
\]
because $\sum_i h_i(t)=2\pi$ and $h_i(t)\le CS$.
\end{proof}

We next record the only finite-element statement about the physical
support trace.  The estimate is on its natural support space, not on an
arbitrary ambient target.

\begin{lemma}[Stable physical material transport]\label{PT:lem:v-transport}
Along every comparable path \eqref{PT:eq:comparable} there is a linear
transport $z_1\mapsto z_t\in Z_{h(t)}$, preserving the exact area and
translation quotient, such that, with
\begin{equation}\label{PT:eq:vtilde}
 \widetilde v_t=v_{h(t)}(z_t)\circ F_t,
\end{equation}
interpreted infinitesimally on the common refined grid,
\begin{align}
 \sup_{0\le t\le1}\|v_{h(t)}(z_t)\|_{H^{1/2}}
 &\le C\|v_h(z_1)\|_{H^{1/2}},\label{PT:eq:v-uniform}\\
 \|\dot{\widetilde v}_t\|_{H^{1/2}}
 &\le C\delta(t)\|v_{h(t)}(z_t)\|_{H^{1/2}}.
 \label{PT:eq:vdot}
\end{align}
\end{lemma}

\begin{proof}
We include the discrete endpoint estimate which prevents a hidden
mesh-frequency loss.  On a comparable cyclic mesh define
\begin{equation}\label{PT:eq:discrete-H}
 \|z\|_{\mathfrak h^{1/2}_S}^2
 =S\sum_i|z_i|^2
 +\sum_{i\ne j}\frac{|z_i-z_j|^2}{(1+|i-j|_{\rm cyc})^2}.
\end{equation}

First let $I_hz$ be the affine interpolant with nodal values $z_i$.
For non-neighbouring cells $I_i,I_j$ at cyclic distance $r$,
comparability gives
\[
 |x-y|\asymp S(1+r),\qquad |I_i|\asymp|I_j|\asymp S.
\]
The corresponding Gagliardo rectangle is therefore comparable with
$|\bar z_i-\bar z_j|^2/(1+r)^2$, where $\bar z_i$ is the cell average.
The two-cell rectangles give $|z_{i+1}-z_i|^2$.  Since
\[
 |\bar z_i-z_i|\le\tfrac12|z_{i+1}-z_i|,
\]
the triangle inequality converts nodal differences to average
differences and conversely, with an error paid by these neighbouring
rectangles.  The summability of $(1+r)^{-2}$ pays each such error a
bounded number of times.
Together with
$\|I_hz\|_2^2\asymp S\sum_i|z_i|^2$, summing the rectangles in both
directions proves
\begin{equation}\label{PT:eq:affine-discrete}
 \|I_hz\|_{H^{1/2}}^2\asymp
 \|z\|_{\mathfrak h^{1/2}_S}^2.
\end{equation}

Now pull every patch $K_i$ in \eqref{PT:eq:trace} to $[0,1]$.  With
$w_i=(h_{i-1}+h_i)/2$, $p_i=h_{i-1}/(h_{i-1}+h_i)$, and
\[
 Z=z_i,\qquad A=w_iA_i^-,\qquad B=w_iA_i^+,
\]
the normalized trace is
\[
 Z\cos\xi+\{(1-y)A+yB\}\frac{\sin\xi}{w_i},
 \qquad \xi=w_i(y-p_i),
\]
where replacing $t_i$ by $y$ changes the coefficient functions by
$O(S^2)$ in $C^1$.  Comparability places $p_i$ in a fixed compact
subinterval of $(0,1)$.  As $w_i\to0$, the local three-node map tends to
\begin{equation}\label{PT:eq:limiting-local-trace}
 (z_{i-1},Z,z_{i+1})\longmapsto
 Z+(y-p_i)\left\{
 (1-y)\frac{Z-z_{i-1}}{2p_i}
 +y\frac{z_{i+1}-Z}{2(1-p_i)}\right\}.
\end{equation}
This limiting map is injective: evaluation at $y=p_i$ first gives
$Z=0$, and division by $y-p_i$ then gives an affine function whose two
endpoint coefficients force $z_{i-1}=z_{i+1}=0$.  Finite-dimensional
norm equivalence, compactness in $p_i$, and an $O(S^2)$ perturbation
therefore give uniform upper and lower bounds for the local map and its
inverse.  The physical trace and the affine interpolant are consequently
related by uniformly invertible three-point finite-element maps.
Applying the preceding rectangle argument to these local maps gives
\begin{equation}\label{PT:eq:discrete-equivalence}
 \|v_h(z)\|_{H^{1/2}}^2
 \asymp\|z\|_{\mathfrak h^{1/2}_S}^2
\end{equation}
on the exact quotient.  Indeed, the area row controls the constant mode,
while support translations are exactly the first two harmonics and are
removed by $\PiNG$.  Hence no low-mode kernel remains in either direction.

We next differentiate the actual coefficients.  Put
\[
 b_i^- =\frac{h_{i-1}}2,\quad b_i^+=\frac{h_i}2,\quad
 D_i=\tan b_i^-+\tan b_i^+,\quad
 N_i(x)=\tan x+\tan b_i^-.
\]
For the pulled local coordinate $x=x(t)$, direct differentiation gives
\begin{align}
 \dot t_i(x)
 &=\frac{\{\sec^2x\,\dot x+\tfrac12\sec^2b_i^-d_{i-1}\}D_i
       -N_i(x)\{\tfrac12\sec^2b_i^-d_{i-1}
                  +\tfrac12\sec^2b_i^+d_i\}}
       {D_i^2},\label{PT:eq:tdot}\\
 \dot A_i^-
 &=\frac{\dot z_i-\dot z_{i-1}}{\sin h_{i-1}}
   -\frac{(z_i-z_{i-1})\cos h_{i-1}}{\sin^2h_{i-1}}d_{i-1}\\
 &\hspace{12mm}-\dot z_i\tan\frac{h_{i-1}}2
   -\frac12z_i\sec^2\frac{h_{i-1}}2d_{i-1},
 \label{PT:eq:Adot-minus}\\
 \dot A_i^+
 &=\frac{\dot z_{i+1}-\dot z_i}{\sin h_i}
   -\frac{(z_{i+1}-z_i)\cos h_i}{\sin^2h_i}d_i\\
 &\hspace{12mm}+\dot z_i\tan\frac{h_i}2
   +\frac12z_i\sec^2\frac{h_i}2d_i.
 \label{PT:eq:Adot-plus}
\end{align}
Moreover
\begin{align}
 \partial_t\cos x&=-\dot x\sin x,\qquad
 \partial_t\sin x=\dot x\cos x,\label{PT:eq:frame-dots}\\
 |\dot x|&\le C\delta(t)S,\qquad
 |\dot t_i|\le C\delta(t).
\end{align}
Substitution of \eqref{PT:eq:tdot}--\eqref{PT:eq:frame-dots} in
\eqref{PT:eq:trace} shows that every pulled trace derivative is a bounded
three-point finite-element combination of $\dot z$ and
\begin{equation}\label{PT:eq:transport-ledger}
 S a_i z_i,\qquad
 a_i(z_{i+1}-z_i),\qquad
 a_{i-1}(z_i-z_{i-1}),
 \qquad |a_i|\le\delta(t),
\end{equation}
and their one-cell shifts.  There is no unweighted term $a_i z_i$: the
trigonometric frame contributes $d_i=S a_i$, while every derivative of a
normalized slope is paired with a nodal difference.  The two one-sided
wall terms cancel because the physical affine vector field is continuous.
Differentiating $\PiNG$ adds only the modes $1,\cos\theta,\sin\theta$ and
their pulled derivatives $u_t\sin\theta,u_t\cos\theta$.  After subtracting
the irrelevant common rotation, periodic Poincar\'e gives
$\|u_t\|_\infty\le C\|u_t'\|_\infty$; these fixed-dimensional rows obey
the same bound as \eqref{PT:eq:transport-ledger} and create no factor of $N$.

The multiplier bound can be proved without a Fourier endpoint theorem.
Write $Dz_i=z_{i+1}-z_i$ and let $K(i,j)=(1+|i-j|_{\rm cyc})^{-2}$.
For $w_i=Sa_iz_i$,
\begin{align*}
 S\sum_i|w_i|^2&\le C\|a\|_\infty^2S\sum_i|z_i|^2,\\
 \sum_{i\ne j}K(i,j)|w_i-w_j|^2
 &\le2S^2\|a\|_\infty^2
      \sum_{i\ne j}K(i,j)|z_i-z_j|^2\\
 &\quad+8S^2\|a\|_\infty^2
      \sum_j|z_j|^2\sum_iK(i,j).
\end{align*}
Since $\sup_j\sum_iK(i,j)<\infty$ and $S$ is bounded, this is controlled
by \eqref{PT:eq:discrete-H}.  For $w_i=a_iDz_i$ use
\[
 S\sum_i|w_i|^2
 \le S\|a\|_\infty^2\sum_i|Dz_i|^2
 \le C\|a\|_\infty^2\|z\|_{\mathfrak h^{1/2}_S}^2
\]
and
\[
 |w_i-w_j|^2
 \le2\|a\|_\infty^2|Dz_i-Dz_j|^2
     +8\|a\|_\infty^2|Dz_j|^2
\]
and
\[
 |Dz_i-Dz_j|^2
 \le2|z_{i+1}-z_{j+1}|^2+2|z_i-z_j|^2.
\]
The last row is controlled by the $|i-j|_{\rm cyc}=1$ part of
\eqref{PT:eq:discrete-H}.  These estimates prove
\begin{equation}\label{PT:eq:ledger-bound}
 \|Sa z+a\Delta z\|_{\mathfrak h^{1/2}_S}
 \le C\|a\|_{\ell^\infty}
       \|z\|_{\mathfrak h^{1/2}_S}.
\end{equation}

It remains to define, rather than merely invoke, the constraint
transport.  Let
\begin{equation}\label{PT:eq:area-row}
 L_tz=\sum_i\ell_i(t)z_i,\qquad
 \ell_i(t)=\tan\frac{h_{i-1}(t)}2+\tan\frac{h_i(t)}2,
\end{equation}
and let
\[
 \tau_t^1=(\cos\theta_i(t))_i,\qquad
 \tau_t^2=(\sin\theta_i(t))_i.
\]
These are the exact two support-translation vectors and satisfy
$L_t\tau_t^\mu=0$.  The matrix
\begin{equation}\label{PT:eq:translation-Gram}
 M_t^{\mu\nu}=\sum_i\ell_i(t)\tau_{t,i}^\mu\tau_{t,i}^\nu
\end{equation}
is uniformly comparable with the identity.  Define
\begin{align}
 \mathcal Q_tz
 &=z-\sum_{\mu,\nu=1}^2(M_t^{-1})_{\mu\nu}
 \left(\sum_i\ell_i(t)z_i\tau_{t,i}^\nu\right)\tau_t^\mu,
 \label{PT:eq:translation-projection}\\
 C_tz
 &=\mathcal Q_t\left(z-\frac{L_tz}{L_t\mathbf1}\mathbf1\right).
 \label{PT:eq:constraint-transport}
\end{align}
If $z_1$ is in the chosen exact gauge, set
\begin{equation}\label{PT:eq:z-transport}
 z_t=C_tz_1.
\end{equation}
Then $L_tz_t=0$, $z_t$ is orthogonal to both translation rows in the
$\ell(t)$-weighted Gram product, and $C_1z_1=z_1$.

The explicit rows give
\begin{align}
 \ell_i&\asymp S,\qquad
 |\dot\ell_i|\le C\delta(t)S,\label{PT:eq:constraint-row-bounds}\\
 \|\dot\tau_t^\mu\|_{\mathfrak h^{1/2}_S}
 &\le C\delta(t),\qquad
 |\dot M_t|+|\partial_tM_t^{-1}|\le C\delta(t).
\end{align}
Indeed, $|\dot\theta_i|=|\sum_{k<i}d_k|\le C\delta(t)$ and adjacent
differences of $\dot\theta_i$ equal $d_i$.  Differentiating
\eqref{PT:eq:translation-projection}--\eqref{PT:eq:constraint-transport}, and
using the bounded area and translation functionals furnished by
\eqref{PT:eq:discrete-equivalence}, first yields
\begin{equation}\label{PT:eq:Ct-raw-bounds}
 \|C_t\|\le C,
 \qquad
 \|\dot C_tz_1\|_{\mathfrak h^{1/2}_S}
 \le C\delta(t)\|z_1\|_{\mathfrak h^{1/2}_S}.
\end{equation}
On $Z_{h(1)}$ the terminal map $C_1$ is the identity.  Hence
\[
 \|(C_t-C_1)|_{Z_{h(1)}}\|
 \le C\int_t^1\delta(s)\,ds
 \le CS^{1/8}.
\]
For sufficiently small $S$ this norm is at most $1/2$.  The Neumann
series therefore proves that $C_t|_{Z_{h(1)}}$ is injective onto its
image and that its inverse is uniformly bounded.  Combining this inverse
bound with \eqref{PT:eq:Ct-raw-bounds} gives
\begin{equation}\label{PT:eq:Ct-bounds}
 \|C_t\|+\|(C_t|_{Z_{h(1)}})^{-1}\|\le C,\qquad
 \|\dot C_tz_1\|_{\mathfrak h^{1/2}_S}
 \le C\delta(t)\|C_tz_1\|_{\mathfrak h^{1/2}_S}.
\end{equation}

Equations \eqref{PT:eq:discrete-equivalence},
\eqref{PT:eq:ledger-bound}, and \eqref{PT:eq:Ct-bounds} now give
\[
 \|\dot{\widetilde v}_t\|_{H^{1/2}}
 \le C\delta(t)\|v_{h(t)}(z_t)\|_{H^{1/2}}.
\]
The finite maps $F_t$ and $F_t^{-1}$ have uniformly bounded Lipschitz
constants by \eqref{PT:eq:comparable}, so their pullbacks are uniformly
bounded on $H^{1/2}$.  Finally,
\begin{equation}\label{PT:eq:delta-integral}
 \int_0^1\delta(t)\,dt
 \le C\frac{\max_i|d_i|}{S}\le CS^{1/8}
\end{equation}
in the near regime.  Gronwall's inequality proves
\eqref{PT:eq:v-uniform}, and the displayed derivative estimate is
\eqref{PT:eq:vdot}.
\end{proof}

\subsection{The near material estimate}

\begin{proposition}[Comparable Lipschitz material bound]
\label{PT:prop:material}
Along the near path \eqref{PT:eq:path} and the transport of
Lemma~\ref{PT:lem:v-transport},
\begin{equation}\label{PT:eq:material}
 \boxed{
 \left|\frac d{dt}
 \Cnull(q_{h(t)},q_{h(t)},v_{h(t)}(z_t))\right|
 \le CS^{5/2}\delta(t)\Gtilde_h(z_1)^{1/2}.}
\end{equation}
\end{proposition}

\begin{proof}
Pull all three functions by the common refined-grid map.  The derivative
splits into the common-coordinate row and three profile rows.
Lemma~\ref{PT:lem:covariance}, \eqref{PT:eq:q-size},
\eqref{PT:eq:uprime}, and \eqref{PT:eq:v-uniform} bound the coordinate row by
the right side of \eqref{PT:eq:material}.

For either fan-profile row use the polarized estimate
\eqref{PT:eq:one-slot}, placing the Lipschitz derivative on the unchanged
fan slot, and then use \eqref{PT:eq:qdot}:
\[
 |\Cnull(q_t,\dot q_t,v_t)|
 \le C S\,(\delta(t)S^{3/2})\|v_t\|_{H^{1/2}}.
\]
For the physical-profile row, use \eqref{PT:eq:one-slot},
\eqref{PT:eq:q-size}, and \eqref{PT:eq:vdot}:
\[
 |\Cnull(q_t,q_t,\dot{\widetilde v}_t)|
 \le C S\,S^{3/2}\delta(t)\|v_t\|_{H^{1/2}}.
\]
Metric equivalence \eqref{PT:eq:metric} completes the proof.
\end{proof}

\begin{lemma}[Regular quotient cancellation]\label{PT:lem:regular}
For the uniform fan and every exact support class,
\begin{equation}\label{PT:eq:regular}
 \Cnull(q_{a\mathbf1},q_{a\mathbf1},v_{a\mathbf1}(z))=0.
\end{equation}
\end{lemma}

\begin{proof}
The fan spline and the physical support construction commute with the
cyclic shift.  Hence the left side is a cyclic-invariant linear functional
of the support coefficients, and therefore a multiple of $\sum_i z_i$.
At the regular tangential polygon the exact area tangent condition is
$\sum_i z_i=0$; translations have already been removed.
\end{proof}

\begin{theorem}[Unconditional mixed row DM]\label{PT:thm:DM}
For every sufficiently fine positive fan and every exact support class,
\begin{equation}\label{PT:eq:DM}
 \boxed{
 |\Cnull(q_h,q_h,v_h(z))|
 \le C(S^{1/4}+S^{1/2})\Jfour(h)^{1/2}
       \Gtilde_h(z)^{1/2}.}
\end{equation}
In particular \eqref{PT:eq:DM-target} holds independently of every minimum
cell and adjacent ratio.
\end{theorem}

\begin{proof}
Use \eqref{PT:eq:far-bound} in the far regime.  In the near regime transport
the final support class to the regular quotient and integrate
\eqref{PT:eq:material}.  Equations \eqref{PT:eq:comparable} and
\eqref{PT:eq:Jidentity} give
\[
 \Jfour(h)^{1/2}\ge cS\left(\sum_i d_i^2\right)^{1/2}
 \ge cS\max_i|d_i|.
\]
Therefore
\begin{align*}
 |\Cnull(q_h,q_h,v_h(z))|
 &\le CS^{5/2}\frac{\max_i|d_i|}{S}
       \Gtilde_h(z)^{1/2}\\
 &\le CS^{1/2}\Jfour(h)^{1/2}\Gtilde_h(z)^{1/2},
\end{align*}
where Lemma~\ref{PT:lem:regular} supplies the zero initial value.  Combining
the regimes proves \eqref{PT:eq:DM}.
\end{proof}

\begin{corollary}[Handoff to the current NR assembly]\label{PT:cor:J0}
Theorem~\ref{PT:thm:DM} supplies the mixed terminal input DM.  The exact flux
audit shows, however, that DM does not imply the former JE--FM clause,
which remains unavailable; its former period-three material falsification
is retracted after the polygonal corner finite-part audit.  NR and active-chart J0 are
assembled conditionally in \texttt{JC\_joint\_comparison.tex} and
\texttt{NR\_assembly.tex}.  At the current audit DF, the CT core, and
JP--DE are closed, while JC--PL is the unique open premise and JP--JC is
its dependent exact output.  No historical CT
proposition is imported in that assembly, and none of these conclusions
uses a minimum fan, polar, or conformal cell.
\end{corollary}
\endgroup
